\section{Introduction}\label{intro:section}

A fibration carries two sources of pluricanonical forms: the geometry
of its base and the variation of its fibers. Iitaka's subadditivity
problem asks for the contribution of the first. The Iitaka--Viehweg
refinement asks how many further directions are forced by a family
whose fibers change birationally. On an open base the correct measure
of the base is its logarithmic Kodaira dimension. We prove this
refinement for projective fibrations over such bases.

For a rational divisor $L$, write $\kappa(L)$ for the Iitaka dimension
of its sufficiently divisible integral multiples. For a smooth
quasi-projective variety $U$, choose a smooth projective compactification
$X$ with reduced simple normal crossing boundary $D_X=X\setminus U$.
A rational boundary is an effective rational divisor with coefficients
in $[0,1]$; SNC abbreviates simple normal crossing support. The
logarithmic Kodaira dimension is
\[
             \kbar(U)=\kappa(X,K_X+D_X);
\]
it is independent of this compactification. We give a point Kodaira
dimension zero and use $(-\infty)+a=-\infty$, including
$a=-\infty$.

Let $f:U\to V$ be a projective surjective morphism with connected
fibers between smooth connected complex quasi-projective varieties.
Fix an algebraic closure $\Omega$ of $\C(V)$, and let $F$ be the smooth
projective geometric generic fiber. Its \emph{birational variation} is
\begin{equation}\label{eq:variation-definition}
 \Var(f)=\min_L\operatorname{trdeg}_{\C}L,
\end{equation}
where $\C\subseteq L\subseteq\Omega$ is algebraically closed and
$F$ is birational to the base extension of a variety over $L$.
This is the field-of-definition form of the classical invariant
\cite{ViehwegAdditivity,KawamataKodaira}. It records the entire function field of $F$;
finite changes of the original base do not change it.

\begin{theorem}[Logarithmic variation]\label{thm:main}
Let $f:U\to V$ be a projective surjective morphism with connected fibers
between smooth connected quasi-projective complex varieties.
Assume $\kbar(V)\geq0$, and let $F$ be its geometric generic fiber.
Then
\begin{equation}\label{eq:main}
 \kbar(U)\geq\kappa(F)+\max\{\kbar(V),\Var(f)\}.
\end{equation}
\end{theorem}

This proves Popa's Conjecture~3.8 in the formulation on page~6 of
his author version dated 13 June 2023 \cite{PopaConjectures}.
The morphism may have singular fibers. The theorem imposes no abundance,
good minimal model, or semiampleness hypothesis on $F$. It concerns the
lower bound in that conjecture; the separate additivity questions for
smooth morphisms in \cite[Conjectures~3.1 and~3.9]{PopaConjectures}
have a different conclusion.
The companion \emph{The reverse logarithmic Kodaira inequality and
additivity} \cite{AdditivityCompanion} proves the corresponding
logarithmic additivity for projective complex reduced-SNC pairs whose total space
and every boundary stratum are smooth over the open base. That
setting allows horizontal boundary and all signs of the base and
fiber Kodaira dimensions. The proof of the variation inequality
here uses the subadditivity and period inputs described below.

The base hypothesis is genuinely logarithmic. For example,
$\kbar(\C^*)=0$, whereas its smooth projective compactification
$\mathbf P^1$ has Kodaira dimension $-\infty$. A compactification
argument must therefore retain the coefficient-one inverse boundary;
replacing the open base by its ordinary projective Kodaira dimension
would lose the hypothesis used in the proof.

\begin{corollary}[Projective Iitaka--Viehweg inequality]
\label{thm:projective-variation}
Let $f:X\to Y$ be a surjective morphism with connected fibers between
smooth connected projective complex varieties, with $\kappa(Y)\geq0$.
For its geometric generic fiber $F$,
\[
 \kappa(X)\geq\kappa(F)+\max\{\kappa(Y),\Var(f)\}.
\]
\end{corollary}

\begin{proof}
Apply Theorem~\ref{thm:main} with empty compactification boundaries.
\end{proof}

Thus the ordinary projective-complex Iitaka--Viehweg $C^+$ conjecture
also holds in every dimension without a good minimal model hypothesis.

\subsection{History and the two constancy problems}

Iitaka developed the divisor-dimension and logarithmic frameworks
for the addition problem \cite{IitakaOriginal,IitakaLog}.
Viehweg's weak positivity of pluricanonical direct images made
birational variation part of the expected lower bound
\cite{ViehwegAdditivity,Viehweg83b}. Koll\'ar proved the variation
refinement for fibers of general type, and Kawamata established the
case of fibers admitting good minimal models
\cite{Kollar87,KawamataKodaira}. These results distinguish a changing
family from a product even when their fiber dimensions agree.

Analytic positivity and generic vanishing brought substantial progress
for special bases. Cao--P\u aun proved subadditivity over abelian
varieties, and Hacon--Popa--Schnell treated bases of maximal Albanese
dimension \cite{CaoPaunAbelian,HaconPopaSchnellAbelian}.
Meng's determinant-based equalities give related refinements over
abelian and maximal-Albanese-dimension bases; the comparison with
birational variation there still uses the good-minimal-model condition
\cite[Theorems~1.9--1.10 and the preceding discussion]{MengSurjective}.
For reduced SNC pairs, Hashizume proved logarithmic subadditivity when
the log canonical divisor of the general fiber is abundant
\cite[Theorem~1.2]{Hashizume}. Popa's formulation asks for the variation
bound over an arbitrary base of nonnegative logarithmic Kodaira dimension
\cite[Conjecture~3.8]{PopaConjectures}.

Cyclic covers and Hodge theory also relate pluricanonical positivity
to variation. Viehweg--Zuo construct Higgs bundles from cyclic covers
associated with pluricanonical sections \cite[\S4]{ViehwegZuo02},
and Popa--Schnell develop this construction using Hodge modules
\cite[\S\S9--11]{PopaSchnellHyperbolicity}. For smooth projective families of
maximal variation whose geometric generic fiber has a good minimal
model, Popa--Schnell prove that the base is of log general type
\cite[Theorem~A]{PopaSchnellHyperbolicity}.

There are also older public claims in the ordinary setting. Tsuji
states full ordinary subadditivity and stronger direct-image generation
in \cite[Theorems~1.9 and~1.13]{TsujiDirectImages}; Maehara states an
ordinary determinant bound involving whole-fiber variation in
\cite[\S7, Theorem~8]{MaeharaIitaka}. These are preprint claims, and their
proofs are not inputs here. The statements just cited do not formulate
the reduced-SNC logarithmic-base inequality of Theorem~\ref{thm:main}.

The proof below starts from the reduced-SNC subadditivity theorem and
the full-period adjoint comparison proved in the companion
\emph{Orbifold and logarithmic Iitaka subadditivity}
\cite{SubadditivityCompanion}. Both interfaces are stated precisely
below. Subadditivity accounts for $\kbar(V)$, but the variation term
requires two additional constancy arguments. The first makes a
Kodaira-dimension-zero fiber constant when its entire top Hodge line
is flat. The second descends the whole original fiber, retaining the
coordinates of its relative Iitaka base. Constancy of either the
canonical model or the Kodaira-zero fiber alone would not give this
last conclusion.

The proof combines several established methods. Canonical
bundle formulas separate discriminant and moduli contributions
\cite{FujinoMori,AmbroBoundary,AmbroModuli}; the root-cover construction
here keeps the actual relative pluriform orders. Deligne's fixed-part
and finite-character results, Schmid's boundary theory, and functorial
Hodge extensions control those lines after restriction
\cite{Deligne71,Schmid,BFMT}. The birational constancy argument uses the singular-weight extension
theorem of Demailly--Hacon--P\u aun, within the Ohsawa--Takegoshi
extension method, and the relative pluricanonical Bergman metrics of
P\u aun--Takayama \cite{DHP13,PaunTakayama18}. The latter develop
the constructions of Berndtsson--P\u aun
\cite{BerndtssonPaunBergman,BerndtssonPaunSubadjunction}. These
estimates are followed by a big klt adjoint model
\cite{BCHM}. For marked model pairs, we use the stable-family positivity
of Kov\'acs--Patakfalvi \cite{KovacsPatakfalvi}. Hanamura's birational
group structure and invariance of finite \'etale covers provide the
final descent mechanism \cite{Hanamura87,Hanamura88,SGA1,Landesman24}.
These inputs are used in the stated settings; the original fiber is
never assumed to have a good minimal model.

\subsection{The proof and its parameter field}

Assume $d=\kappa(F)\geq0$. Choose compatible compactifications
$f:(X,D_X)\to(Y,D_Y)$, retaining $U=f^{-1}(V)$.
Logarithmic subadditivity gives $\kbar(U)\geq d+\kbar(V)$.
Let $q:X'\to Z$ be the absolute Iitaka map of $K_X+D_X$ on a
smooth model, so $\dim Z=\kbar(U)$. Take the relative algebraic closure of $\C(Y)\C(Z)$ in $\C(X)$,
and resolve the normalization of the combined image in $Y\times Z$
in this field. On compatible smooth models this gives
\[
             X'\xrightarrow{x}W,
       \qquad g:W\to Y,\qquad h:W\to Z.
\]
The generic $x$-fiber $J$ has ordinary Kodaira dimension zero. The
images $g(W_z)$ form an algebraic family of subvarieties of $Y$.
Let $T_0$ be a smooth projective model of the normalization of its
image parameter space in the relative algebraic closure of that
parameter field in $\C(Z)$. Contraction of a
nonzero logarithmic pluriform on $Y$ controls the differential of this
family and proves
\begin{equation}\label{intro:parameter-dimension}
       \kbar(U)=d+\dim T_0,
       \qquad \C(T_0)\subseteq\C(Y).
\end{equation}
Section~\ref{it:section} constructs the maps and establishes the
inclusion as an inclusion of the actual function fields.

It remains to define the whole fiber $F$ over
$b=\overline{\C(T_0)}\subset\Omega$. Choose the least $p>0$ with a nonzero section of $pK_J$, adjoin
its $p$th root as a canonical differential, and resolve the resulting
cover. This minimal ordinary root cover has an entire top-form space
of dimension one. The adjoint
comparison, applied anew to the restricted integral variation, makes
its Hodge line rationally trivial along the $h$-fibers. Finite
character then follows on that restricted algebraic locus.

The analytic criterion of Section~\ref{cc:section} turns this flatness
into birational constancy of the root covers, equivariantly for their
deck action. A lifted base vector field may initially have poles.
Minimizing an $L^{2/m}$ integral while fixing a leading coefficient
along a divisor shows that every pole has positive asymptotic fixed
order. A single sufficiently small ample perturbation of the canonical
divisor gives a big adjoint divisor whose model contracts these poles;
the resulting regular vector field has flows identifying nearby fibers.
The argument is carried out with the weight chosen before the final
fiber and with the regularization, degree, and cutoff limits in a
fixed order.

The last step is algebraic. Marked canonical models recover a finite
normalization $V'$ of the relative Iitaka base over $b$, together with
its map and the images of the old boundary. The family $J$ is then a
form of a reference variety over $k=b(V')$. A finite splitting cocycle
can be made constant. At every divisor which moves with the parameter,
the actual relative-form order calculation supplies a multiplicity-one
component attaining the threshold. This forces the splitting cover to
be unramified there. After removing finitely many fixed divisors,
purity and descent of finite \'etale covers produce a field $F_b$ with
\[
 k\subseteq F_b,
 \qquad
 \operatorname{Frac}(\Omega\otimes_bF_b)
   \simeq\operatorname{Frac}(\Omega\otimes_{\C(Y)}\C(X)).
\]
Thus the entire geometric generic fiber is defined birationally over
$b$, and $\Var(f)\leq\dim T_0$. Together with
\eqref{intro:parameter-dimension} and the subadditivity bound, this
proves Theorem~\ref{thm:main}.

Figure~\ref{fig:whole-fiber-route} records the two constancy steps and
the final field inclusion in this argument.

\begin{figure}[htbp]
\centering
\begin{tikzpicture}[
  box/.style={draw,rounded corners=2pt,align=center,text width=.40\textwidth,
              inner sep=6pt,font=\small},
  arrow/.style={-{Stealth[length=2mm]},thick},node distance=12mm and 9mm]
\node[box] (parameter) {Absolute logarithmic Iitaka system\\
 $\kbar(U)=d+\dim T_0$};
\node[box,right=of parameter] (hodge) {Restricted root-cover comparison\\
 Flat top line; adjoint section\\dimension zero};
\node[box,below=of parameter] (marked) {Marked canonical models\\
 Recover $k=b(V')$};
\node[box,below=of hodge] (constant) {Analytic constancy criterion\\
 $J$ is geometrically birationally\\constant on the $h$-fibers};
\node[box,below=of marked] (cover) {Multiplicity-one threshold minimizer\\
 No ramification at moving divisors};
\node[box,below=of constant] (whole) {Finite \'etale descent\\
 $k\subseteq F_b$; the whole $F$ is birationally defined over $b$};
\draw[arrow] (parameter) -- (hodge);
\draw[arrow] (hodge) -- (constant);
\draw[arrow] (parameter) -- (marked);
\draw[arrow] (hodge) -- (marked);
\draw[arrow] (marked) -- (cover);
\draw[arrow] (cover) -- (whole);
\draw[arrow] (constant) -- (whole);
\end{tikzpicture}

\caption{The parameter field and the two constancy steps. The restricted
root-cover comparison supplies a flat line for analytic constancy and
a zero-dimensional adjoint section system for marked-model recovery.
The final field retains the relative Iitaka-base coordinates, giving
a bound for the variation of the whole $F$.}
\label{fig:whole-fiber-route}
\end{figure}

\subsection{Smooth families}

For smooth families, the good-model theorem of the companion
\emph{Log abundance in characteristic zero} and Taji's rigidity
theorems give a further application. This application is not used in
the proof of Theorem~\ref{thm:main}. A smooth quasi-projective base $V$ is
\emph{Campana-special} if, for any smooth projective compactification
$(Y,D)$ with reduced SNC boundary, every invertible subsheaf
$L\subseteq\Omega_Y^p(\log D)$ satisfies $\kappa(L)<p$ for
$1\leq p\leq\dim V$ \cite[\S1]{TajiGoodModels}.

\begin{corollary}[Variation and birational isotriviality]
\label{cor:nonuniruled-rigidity}
Let $f:U\to V$ be a smooth projective surjection with connected fibers
between smooth connected complex quasi-projective varieties. Assume
that every closed fiber is non-uniruled. Then
\[
\begin{aligned}
 \kbar(V)=-\infty&\quad\Longrightarrow\quad \Var(f)<\dim V,\\
 \kbar(V)\geq0&\quad\Longrightarrow\quad \Var(f)\leq\kbar(V).
\end{aligned}
\]
If $V$ is Campana-special, then $\Var(f)=0$. In particular, if
$\kbar(V)=0$, the family is birationally isotrivial.
\end{corollary}

\begin{proof}
Fix a closed fiber $F_v$. By \cite[Corollary~0.3]{BDPP}, its
non-uniruledness makes $K_{F_v}$ pseudo-effective. The companion
\cite[Corollary~11.2]{LogAbundanceCompanion} gives a normal projective
model $M_v$ with semiample $\Q$-Cartier $K_{M_v}$ and, on a common
smooth resolution $p:W\to F_v$, $q:W\to M_v$,
\[
 p^*K_{F_v}=q^*K_{M_v}+E,\qquad E\geq0
 \quad\text{and $E$ is $q$-exceptional}.
\]
For any prime divisor $P$ exceptional over $M_v$, pass to a higher
common resolution containing $P$ and denote the effective pullback
of $E$ again by $E$. Since $F_v$ is smooth, the discrepancy identity
\[
 a(P,M_v)=a(P,F_v)+\operatorname{coeff}_P E\geq0
\]
shows that $M_v$ has canonical singularities. Thus every $F_v$ admits
a good minimal model in the sense used by Taji. His variation is the
same whole-fiber field-of-definition invariant as
\eqref{eq:variation-definition}, so
\cite[Theorems~1.1--1.2]{TajiGoodModels} gives the dichotomy and the
special-base assertion. The case $\kbar(V)=0$ also follows from the
second inequality and $\Var(f)\geq0$.
\end{proof}

Here \emph{birationally isotrivial} means precisely $\Var(f)=0$.
No isomorphism or trivialization of the family, globally or after a
finite \'etale cover, is asserted.

For a smooth family in Theorem~\ref{thm:main} with $\kappa(F)\geq0$,
the nonnegative-base bound also follows directly from logarithmic
additivity \cite[Corollary~1.2]{AdditivityCompanion}. Choose compatible
smooth compactifications with $U=f^{-1}(V)$ and reduced complement
boundaries; their boundary strata are empty over $V$. The companion
gives $\kbar(U)=\kappa(F)+\kbar(V)$, where the geometric generic and
very general fiber Kodaira dimensions agree by generic base change and
upper semicontinuity in each pluricanonical degree. Comparing with
\eqref{eq:main} and cancelling the finite $\kappa(F)$ gives
$\Var(f)\leq\kbar(V)$. This comparison
does not address the negative-base branch or all special bases.

\subsection{Organization and an additional numerical result}

Section~\ref{it:section} constructs the parameter field, after stating
the exact logarithmic lower bound and the section conventions.
Section~\ref{sec:rankone} constructs the least-index root cover,
computes its parabolic line from actual relative orders, and proves
restricted flatness and finite character.
Section~\ref{cc:section} proves the analytic birational constancy
criterion, including contraction, equivariance and algebraic spreading.
Section~\ref{num:comparisons} develops the comparison for auxiliary
big model pairs at its first use. Section~\ref{sec:marking} uses
recognizable markings to recover the finite Iitaka-base normalization.
Section~\ref{ds:section} makes the cocycle constant, removes moving
ramification after every finite parameter extension, and identifies
the whole original function field. This completes the main theorem.

The final two sections give a separate extension. In
Section~\ref{num:second-base}, a nef contribution pulled back from an
adjoint-positive base replaces the Hodge line: interpolation retains
the full second Iitaka-base dimension. Its volume estimate bounds only
the effective action on the section image and is independent of the
degree of the auxiliary alteration. Section~\ref{ordinary:reductions}
applies that theorem to the ordinary relative Iitaka fibration and
recovers the same embedded parameter field. The exact reduced-root
and section-comparison inputs from the companion are stated there.
These sections are additional results; the logarithmic proof above
does not depend on them.
